% Work/MAT, 1. rocnik, hodina 26: soucin, podil, mocnina mocniny a soucinu.
% Revize podle zadani a obrazku z 30. 10. 2026.
% Navazuje na 25_MOC_MocninaSPrirozenymExponentem.tex: pravidla doplnujeme
% jako radky teze prubezne tabulky, bez opakovani definic z hodiny 25.
% Odvozovani: ucebnice/01_Zakladni_poznatky/Mocniny_a_odmocniny/
% Mocniny a Odmocniny.tex; Realisticky.cz: 010602 Vzorce pro mocniny I.
% Vybrane ulohy S2 a D1, priklady P3 a P5 prevzaty ze screenshotu uzivatele.
\documentclass[11pt,a4paper]{article}
% Spolecny styl pro ucitelske pripravy MAT (Work/MAT)
% Verze 1.1, 2026-10-23
\NeedsTeXFormat{LaTeX2e}
\ProvidesPackage{priprava}[2026/10/23 v1.1 Sablona ucitelskych priprav MAT]

% Zakladni typografie a cestina
\RequirePackage[T1]{fontenc}
\RequirePackage[utf8]{inputenc}
\RequirePackage[czech]{babel}
\RequirePackage{lmodern}
\RequirePackage{microtype}
\RequirePackage{amsmath,amssymb}
\RequirePackage{xcolor}
\RequirePackage{enumitem}
\RequirePackage[most]{tcolorbox}
\RequirePackage{tabularx,booktabs,array}
\RequirePackage{fancyhdr}
\RequirePackage{needspace}
\RequirePackage{geometry}
\RequirePackage[hidelinks]{hyperref}

\geometry{a4paper,top=17mm,bottom=18mm,left=23mm,right=23mm}
\definecolor{notegray}{gray}{0.39}
\setlength{\parindent}{0pt}
\setlength{\parskip}{2.5pt plus 1pt minus 0.5pt}
\setlist[itemize]{topsep=2pt,itemsep=1pt,parsep=0pt,leftmargin=1.45em}
\setlist[enumerate]{topsep=2pt,itemsep=2pt,parsep=0pt,leftmargin=1.65em}
\setcounter{secnumdepth}{0}
\renewcommand{\arraystretch}{1.12}

% Zahlavi: vlevo rocnik | globalni tema, vpravo cislo: nazev hodiny
% Zapati zustava zcela prazdne.
\pagestyle{fancy}
\fancyhf{}
\renewcommand{\headrulewidth}{0.35pt}
\renewcommand{\footrulewidth}{0pt}
\setlength{\headheight}{14pt}
\newcommand{\nastavPripravu}[4]{%
  \fancyhead[L]{\small #1\enspace\textbar\enspace #2}%
  \fancyhead[R]{\small #3: #4}%
}

% Minimalisticke cernobile boxy. Kazda definice zvlast.
\tcbset{commonlesson/.style={enhanced,breakable=false,colback=white,colframe=black,
  colbacktitle=white,coltitle=black,boxrule=0.45pt,
  arc=0.8mm,outer arc=0.8mm,left=3mm,right=3mm,top=1.6mm,bottom=1.6mm,
  toptitle=1mm,bottomtitle=0.7mm,fonttitle=\bfseries,
  attach boxed title to top left={xshift=2.5mm,yshift=-1.15mm},
  boxed title style={colback=white,colframe=white,boxrule=0pt,left=1mm,right=1mm}}}
\newtcolorbox{zadaniBox}[1][]{commonlesson,title={Zadání},#1}
\newtcolorbox{prikladBox}[1][]{commonlesson,title={Příklad},#1}
\newtcolorbox{definiceBox}[1][]{commonlesson,title={Definice},#1}

% Spolecny box pro zadani a jeho reseni, bez mezery mezi castmi.
% Pouziti:
%   \begin{prikladReseniBox}
%     Zadani ulohy.
%     \tcblower
%     \begin{align*} ... \end{align*}
%   \end{prikladReseniBox}
% Spodni cast ma jemny sedy podklad a delici linku; lze v ni sazet
% vice radku matematiky. Obe casti maji jen jeden spolecny ramecek.
\tcbset{solutionlesson/.style={
  commonlesson,
  enhanced jigsaw,
  breakable,
  colbacklower=black!3,
  segmentation style={solid,draw=black!25,line width=0.35pt},
  fontlower=\small\color{notegray},
  before lower={\noindent\textbf{Řešení:}\enspace},
  leftlower=3mm,rightlower=3mm
}}
\newtcolorbox{prikladReseniBox}[1][]{solutionlesson,title={Příklad},#1}
\newtcolorbox{zadaniReseniBox}[1][]{solutionlesson,title={Zadání},#1}

% Pomucky pro autora pripravy: sedive poznamky; \reseni pro starsi pripravy
\newcommand{\poznamka}[1]{\par\begingroup\small\color{notegray}\itshape
  \textnormal{\textbf{Poznámka pro učitele.}}\enspace #1\par\endgroup}
\newcommand{\reseni}[1]{\par\begingroup\small\color{notegray}
  \textbf{Očekávané řešení:}\enspace #1\par\endgroup}
\newcommand{\blok}[1]{\needspace{5\baselineskip}\subsection*{#1}}
\newcommand{\vyrok}[1]{\uv{#1}}

\nastavPripravu{1. ročník}{Mocniny a odmocniny}{26}{Pravidla pro mocniny}
\begin{document}

\poznamka{\textbf{Jádro 45 minut:} 3 min aktivace z hodiny 25; 8 min odvození
součinu a podílu + tabulka; 6 min P1 a P2; 11 min S2 (a--c), další dvě
položky jako výzva; 7 min mocnina mocniny + P3 a), b); 7 min mocnina
součinu + P5 a), b); 3 min shrnutí a domácí práce. P3 c), P5 c)--e),
původní S1 a R1 jsou rezerva. Všechny úlohy jsou v přípravě, ale nejde
o povinnou náplň jediné hodiny. V hodině 27 lze pravidla důkladněji procvičit.}

\blok{Navážeme na 25. hodinu: kolik máme činitelů?}
Z minulé hodiny víme, že $a^n$ je součin $n$ stejných činitelů.
Když dva takové součiny spojíme, počty činitelů \textbf{sčítáme}.
Při dělení můžeme společné nenulové činitele odebrat.
\[
\underbrace{a\cdot a\cdot a}_{3}\cdot
\underbrace{a\cdot a}_{2}=a^5,
\qquad
\frac{\overbrace{a\cdot a\cdot a\cdot a\cdot a}^{5}}{\underbrace{a\cdot a}_{2}}=a^3
\quad(a\ne0).
\]
\textbf{Důkaz obecně:} pro $r,s\in\mathbb N$ platí
\[
\underbrace{a\cdots a}_{r\text{ činitelů}}\cdot
\underbrace{a\cdots a}_{s\text{ činitelů}}
=\underbrace{a\cdots a}_{r+s\text{ činitelů}}.
\]
Při podílu $r>s$ a $a\ne0$ zbude po odstranění $s$ stejných
činitelů $r-s$ činitelů v čitateli.

\blok{Nové řádky do naší tabulky pravidel}
Stejně jako v hodině 25 si tabulku \textbf{postupně rozšiřujeme}.
Následující čtyři řádky připojíme k pravidlům pro znaménka,
první mocninu, $0^n$ a $1^n$; nejde o jejich náhradu.
\begin{center}
\small
\begin{tabularx}{\linewidth}{@{}>{\raggedright\arraybackslash}p{31mm}>{\raggedright\arraybackslash}p{50mm}X@{}}
\toprule
\textbf{Pravidlo} & \textbf{Zápis} & \textbf{Podmínky a význam} \\
\midrule
Součin stejného základu & $a^r\cdot a^s=a^{r+s}$ & $a\in\mathbb R$, $r,s\in\mathbb N$ \\
Podíl stejného základu & $\dfrac{a^r}{a^s}=a^{r-s}$ & $a\ne0$, $r>s$, $r,s\in\mathbb N$ \\
Mocnina mocniny & $(a^r)^s=a^{r\cdot s}$ & $a\in\mathbb R$, $r,s\in\mathbb N$ \\
Mocnina součinu & $(ab)^r=a^r b^r$ & $a,b\in\mathbb R$, $r\in\mathbb N$ \\
\bottomrule
\end{tabularx}
\normalsize
\end{center}
\textbf{Případ $r=s$:} pro $a\ne0$ máme $a^r/a^r=1$.
\textbf{Případ $r<s$:} zatím výsledek zapíšeme jako převrácenou
hodnotu mocniny; \emph{záporný exponent} systematicky zavedeme v hodině 28.

\blok{Řešené příklady — zachované z původní přípravy}
\begin{prikladReseniBox}[title={P1: součin a podíl stejných základů},breakable=false]
Uprav $2^5\cdot 2^3$ a $\dfrac{3^7}{3^4}$.
\tcblower
$2^5\cdot2^3=2^{5+3}=2^8=256$;
$\dfrac{3^7}{3^4}=3^{7-4}=3^3=27$.
\end{prikladReseniBox}
\poznamka{\textbf{Otázka žákům právě u P1:} Co kdybychom v podílu
$3^7/3^4$ \emph{prohodili čitatele a jmenovatele}?
Výsledek už nebude $3^3$; dostaneme $3^4/3^7=1/3^3$.
Zatím nech jen vysvětlit počet činitelů. Zápis $3^{-3}$
a jeho pravidla probereme v hodině 28.}

\begin{prikladReseniBox}[title={P2: původní podmínka se neztrácí},breakable=false]
Uprav $\dfrac{x^7x^2}{x^4}$ a uveď podmínku.
\tcblower
Původní jmenovatel vyžaduje $x\ne0$. Potom
$\dfrac{x^7x^2}{x^4}=x^{7+2-4}=x^5$.
I když ve výsledku žádný jmenovatel není, $x\ne0$ nadále platí.
\end{prikladReseniBox}

\blok{S2: výrazy z pracovního listu — povinné a dobrovolné}
\begin{zadaniReseniBox}[title={S2: Odstraň zlomky pomocí pravidel pro mocniny},breakable=false]
Odstraň zlomky ve výrazech (nepoužívej přímé krácení,
ale \textbf{vzorec pro podíl mocnin}). Výsledek udej jako mocninu
nebo součin mocnin. Nezapomeň na podmínky.
\begin{enumerate}[label=\alph*)]
\item $\displaystyle\frac{2^{14}}{2^6}$
\item $\displaystyle\frac{(-4)^{35}}{(-4)^{13}}$
\item $\displaystyle\frac{a^3\cdot b^5}{b^2\cdot a}$
\item $\displaystyle\frac{x^4\cdot(-x)^7}{(-x)^3\cdot x^2}$
\item $\displaystyle\frac{a^4\cdot(-a)^3\cdot b^5}{(-a)^6\cdot(-b)^3}$
\end{enumerate}
\tcblower
\begin{enumerate}[label=\alph*),itemsep=2pt]
\item $2^{14-6}=\boxed{2^8}$.
\item $(-4)^{35-13}=\boxed{(-4)^{22}}$.
\item $a^{3-1}b^{5-2}=\boxed{a^2b^3}$; $a,b\ne0$.
\item $x^{4-2}(-x)^{7-3}=x^2(-x)^4=\boxed{x^6}$; $x\ne0$.
\item $a^4\,(-a)^{3-6}\,b^5\,(-b)^{-3}$ takto \emph{zatím nezapisujeme}.
Znaménka oddělíme ještě před podílem: nahoře $-a^7b^5$,
dole $-a^6b^3$, proto $\boxed{ab^2}$; $a,b\ne0$.
\end{enumerate}
\end{zadaniReseniBox}
\poznamka{V povinném jádru S2 a)--c). Písmena d), e) jsou náročnější
příklady pro rychlé; u e) záměrně nepiš záporné exponenty, ale přeuspořádej
znaménka a spočti rozdíly kladných exponentů. U všech podílů
původně nesmí být jmenovatel nula.}

\blok{Mocnina mocniny — další rozšíření tabulky}
V zápisu $(a^r)^s$ opakujeme součin $a^r$ právě $s$-krát.
V každém bloku je $r$ činitelů, dohromady tedy $r\cdot s$.
\[
(a^r)^s=\underbrace{a^r\cdot a^r\cdots a^r}_{s\text{ bloků}}
=\underbrace{a\cdot a\cdots a}_{r\cdot s\text{ činitelů}}
=a^{r\cdot s}.
\]
\textbf{Pozor:} exponenty zde \emph{násobíme}, nesčítáme.

\begin{prikladReseniBox}[title={P3: Odstraň závorky, výsledek jako mocninu},breakable=false]
\begin{enumerate}[label=\alph*)]
\item $(2^3)^5$
\item $\bigl((-\pi)^3\bigr)^4$
\item $\bigl(-(3)^4\bigr)^3$
\end{enumerate}
\tcblower
\begin{enumerate}[label=\alph*),itemsep=2pt]
\item $(2^3)^5=2^{3\cdot5}=\boxed{2^{15}}$.
\item $\bigl((-\pi)^3\bigr)^4=(-\pi)^{3\cdot4}
=(-\pi)^{12}=\boxed{\pi^{12}}$.
\item $\bigl(-(3)^4\bigr)^3=(-3^4)^3=-\bigl(3^4\bigr)^3
=\boxed{-3^{12}}$. Mínus není součástí základu v $3^4$;
vnější exponent $3$ je lichý.
\end{enumerate}
\end{prikladReseniBox}

\blok{Mocnina součinu — odvoď pravidlo, pak použij}
Pro $r\in\mathbb N$ máme $r$ kopií stejného součinu $(ab)$.
Díky komutativitě a asociativitě násobení přeuspořádáme činitele:
\[
(ab)^r=\underbrace{(ab)(ab)\cdots(ab)}_{r\text{ činitelů }ab}
=\underbrace{a\cdots a}_{r\text{ činitelů}}
 \underbrace{b\cdots b}_{r\text{ činitelů}}=a^r b^r.
\]
\textbf{Důležité:} toto pravidlo funguje pro \emph{součin}, nikoliv
pro součet: $(a+b)^r\ne a^r+b^r$ obecně (např. pro $a=b=1$, $r=2$).

\begin{prikladReseniBox}[title={P5: Mocnina součinu — příklady ze zadání},breakable=false]
Vypočti nebo uprav; výsledek udej jako součin mocnin.
\begin{enumerate}[label=\alph*)]
\item $(2\cdot3)^2$
\item $(2^3\cdot3)^4$
\item $3(a^2\cdot b^3)^3$
\item $\displaystyle\frac{(a^2\cdot b^3)^3\cdot a^4}{(a^3\cdot b^2)^3}$
\item $[(-a)\cdot b]^3(-a^4b^3)^5$
\end{enumerate}
\tcblower
\begin{enumerate}[label=\alph*),itemsep=2pt]
\item $(2\cdot3)^2=\boxed{2^2\cdot3^2}=36$.
\item $(2^3\cdot3)^4=(2^3)^4\cdot3^4
=\boxed{2^{12}\cdot3^4}$.
\item $3(a^2b^3)^3=3(a^2)^3(b^3)^3=\boxed{3a^6b^9}$.
\item Pro $a,b\ne0$:
$\displaystyle\frac{a^6b^9a^4}{a^9b^6}
=a^{6+4-9}b^{9-6}=\boxed{ab^3}$.
\item $[(-a)b]^3(-a^4b^3)^5
=(-a)^3b^3\,(-1)^5a^{20}b^{15}
=\boxed{a^{23}b^{18}}$.
\end{enumerate}
\end{prikladReseniBox}
\poznamka{P5 a), b) patří do jádra; c)--e) jsou výzva
nebo rezerva do hodiny 27. Zvlášť u e) rozlišuj $(-a)^4$ a $-a^4$.
Pravidlo pro mocninu \emph{podílu} zatím nemusíme zavádět.}

\blok{Další původní procvičení — dobrovolná rezerva}
\begin{zadaniReseniBox}[title={S1: krátké základní úpravy z původní přípravy},breakable=false]
\begin{enumerate}[label=\alph*)]
\item $a^3a^5$.
\item $\dfrac{5^6}{5^2}$ a $\dfrac{2^3}{2^5}$.
\item $\dfrac{b^4b^3}{b^5}$ a $\dfrac{b^2}{b^2}$.
\end{enumerate}
\tcblower
\begin{enumerate}[label=\alph*)]
\item $a^8$ pro libovolné $a\in\mathbb R$.
\item $5^4=625$ a $1/2^2=1/4$.
\item $b^2$ a $1$, vždy při $b\ne0$.
\end{enumerate}
\end{zadaniReseniBox}

\begin{zadaniReseniBox}[title={R1: další výrazy z původní přípravy},breakable=false]
\begin{enumerate}[label=\alph*)]
\item $a^4a^3a^2$, $\dfrac{x^9}{x^4}$ a $\dfrac{y^3}{y^7}$.
\item $\dfrac{3^5\cdot3^2}{3^4}$ a $\dfrac{2^3\cdot4^2}{2^5}$.
\item Platí $a^3+a^2=a^5$ pro každé $a$?
\item $\dfrac{a^6b^3\cdot ab^2}{a^4b^6}$; uveď podmínky.
\item Porovnej $\dfrac{7^8}{7^5}$ a $7^2$.
\end{enumerate}
\tcblower
\begin{enumerate}[label=\alph*)]
\item $a^9$, $x^5$ pro $x\ne0$, $1/y^4$ pro $y\ne0$.
\item $3^3=27$ a $2^2=4$.
\item Ne, například pro $a=2$ je $8+4=12\ne32$.
\item $a^3/b$, kde $a\ne0$ a $b\ne0$.
\item $7^3>7^2$.
\end{enumerate}
\end{zadaniReseniBox}

\blok{Domácí procvičení — přesně čtyři položky ze zadání}
\begin{zadaniReseniBox}[title={D1: Zjednoduš výrazy},breakable=false]
\begin{enumerate}[label=\alph*)]
\item $\displaystyle\frac{2^5\cdot2^7}{2^{10}}$
\item $\displaystyle\frac{(-3)^3\cdot(-3)^6}{(-3)^5\cdot3^2}$
\item $\displaystyle\frac{3^{15}}{3^2\cdot3^7}$
\item $\displaystyle\frac{(-2)^5\cdot2^6}{(-2)^4\cdot2^2}$
\end{enumerate}
\tcblower
\begin{enumerate}[label=\alph*),itemsep=2pt]
\item $2^{5+7-10}=\boxed{2^2}=4$.
\item $\dfrac{(-3)^{3+6-5}}{3^2}=\dfrac{(-3)^4}{3^2}
=\boxed{3^2}=9$.
\item $3^{15-2-7}=\boxed{3^6}$.
\item $\dfrac{(-2)^5\cdot2^6}{(-2)^4\cdot2^2}
=(-2)^{5-4}2^{6-2}=(-2)\cdot2^4=\boxed{-2^5}=-32$.
\end{enumerate}
\end{zadaniReseniBox}
\end{document}
